\chapter{Conclusion}

Chunk size produced no detectable difference and retrieval strategy none without
enrichment, while indexing representation moved retrieval substantially, unevenly
across strategies, and with losers as well as gains: some questions retrieved
under text-only were lost under enrichment.

The nulls for chunk size and retrieval strategy have one cause: across
near-identical calls, the company and the quarter reached the retrievers only
where a transcript's own words named them, and neither choice, as varied here,
could supply them.

Retrieval improved substantially with enrichment, but no difference in the
reported arm's answer outcomes between conditions was detectable.

What an evaluation rewards has to be read against what was retrieved. When the
context cannot support an answer, abstaining is the correct response, and for a
conjunctive question that holds even when part of its evidence was retrieved.
Correct abstention is a target to design an evaluation around, not a shortfall
to minimize, and a scheme that counts it as a failure misjudges the system it
scores. An automatic grounding check is likewise an instrument with a direction
of error: a substantial share of what it called ungrounded was faithful text it
could not read.

For teams retrieving over documents that differ mainly in entity and period, the
place to start is the identity metadata they already hold, put into the string
the retriever indexes. The gain to expect falls on questions that name one
company; where a question names none, nothing was detectable here. A
retrieval-augmented system's answers are bounded by what
its retriever finds, and its retriever by what it is given to index.
